\chapter{Lean 4}
\label{c:chapter4}
\todo{Fragen, auf welche ich eine Antwort wissen sollte}: 1) warum sollten/können wir Lean vertrauen.
2) auf welchen mathematischen Grundlagen basiert Lean?3) wie sind die Natürlichen zahlen definiert


Some knowledge you should mention:
\begin{itemize}
  \item Lean compiles to C
  \item the build tool is \texttt{lake} and allows metaprogramming
\end{itemize}

The main purpose of using Lean these days is as proof assistant, but Lean is
also a functional programming language, where every function must 
provably terminate. (from claude.ai: because one could prove FALSE with a looping
function \todo{TODO} HOW?) As a consequence, one has to prove (or convince) Lean that 
the recursion terminates. This is done by assuring every recursive call is done on a subset 
of the argument, to make sure the recursive calls will end on a base case and 
terminate. For illustration, let's look at the quicksort algorithm.

\begin{lstlisting}[float=tbh, mathescape=true, captionpos=b, caption=Lean 4, label=gamma, language=lean]
partial def qsort : List Nat -> List Nat
  | [] => []
  | p :: xs =>
      qsort (xs.filter (fun x => x ≤ p))
      ++ [p]
      ++ qsort (xs.filter (fun x => x > p))
\end{lstlisting}

With keyword \texttt{partial} we can postpone the proof of termination
and evaluate the function

\begin{lstlisting}[float=tbh, mathescape=true, captionpos=b, caption=Lean 4, label=alpha, language=lean]
#eval qsort [1,3,2] --leading to [1,2,3]
\end{lstlisting}

\todo{Haskell has type class Ord, how does Lean handle 
ordinary types? Do I need for each type a new definition?}

Showing that qsort is a correct sorting algorithm, requires Lean 
to tell, what this means. Thus we have to prove the result is sorted
and a permutation of the input list. 

\subsubsection{Lean for proving programs}
The core principle is, to write program specs as theorems and
prove them. In the case of the quicksort algorithem something like

\begin{lstlisting}[float=tbh, mathescape=true, captionpos=b, caption=Lean 4, label=beta, language=lean]
theorem 
is_sorted (ns: List Nat) : (qsort ns).Pairwise (· ≤ ·) :=...
theorem 
is_permutation ...
\end{lstlisting}
\label{qsort_validation}

\todo{TODO}
An example of a theorem in Lean.


 

\section{Basic Language Concepts}

\subsection{Proof by Contradiction}

\begin{lstlisting}[float=tbh, mathescape=true, captionpos=b, caption=Proof by contradiction, label=contradiction, language=lean]
example (P : Prop) (h : ¬ P -> False) : P := by
  by_contra hP 
  exact h hP        
\end{lstlisting}



\section{Definitions and Computation}

\section{The natual numbers}
How are those defined in lean and what can we do with them. Why are the real number so different.

\todo{add here how the nat numbers are defined in lean. How is the link done
from Number 12 ist $succ (... 12 x (succ 0) ..)$}


The only way to make numbers are $0$ and $succ$.

\subsection{Proving 2 + 2 = 4 an example}
Let's look at an example from the \textit{Lean Number Game} to prove $2+2=4$.

\begin{lstlisting}[float=tbh, mathescape=true, captionpos=b, caption=Proof by contradiction, label=contradiction, language=lean]
example : (2 : ℕ) + 2 = 4 := by
    nth_rewrite 2 [two_eq_succ_one] -- only change the second `2` to `succ 1`.
    rw [add_succ]
    rw [one_eq_succ_zero]
    rw [add_succ, add_zero] -- two rewrites at once
    rw [← three_eq_succ_two] -- change `succ 2` to `3`
    rw [← four_eq_succ_three]
    rfl
\end{lstlisting}

\section{Proof Construction}
Lean provides two proof styles
\begin{itemize}
  \item Term mode: direct writing of the term.
  \item Tactics mode: builds the proof step by step, using tactics:
  intro, apply, exact, cases, induction, rw, simp, rfl, calc
\end{itemize}

\todo{what's with tactics for contradiction?}

(claude.ai) Tactics can be arbitrarily clever, but every proof is
rechecked by a small kernel. 

\section{Automation}
Linear arithmetic is handled by \texttt{omega}. Rewriting using 
lemmas is with \texttt{simp}. \todo{and more}



\section{Why should we trust Lean 4}
To trust such a mechanical verification process, it is enough to show that the verification algorithm actually 
verifies that all formal rules are correctly applied. \cite[Chapter~1.2]{coq_book}

On the lean home page \cite{web:lean_homepage} the following is stated: \textit{ Lean's minimal trusted kernel guarantees absolute correctness in mathematical proof, software and hardware verification.}

What does this mean now?

Ref: lean FAQ.
What is the Trusted Code Base of Lean?

One of Lean's core design approaches is that \textbf{one does not 
need to trust Lean's implementation} in order to trust the proofs 
it produces. Proofs can be verified using independent checkers, 
as described in more detail in the section 
\href{https://lean-lang.org/doc/reference/latest/find/?domain=Verso.Genre.Manual.section\&name=validating-proofs}
{``Validating a Lean Proof''} of the Lean Reference Manual.

Thus, it is sufficient to trust the soundness of the checker used 
to validate the proof output. This basis of trust, i.e., the 
source code of the checker(s) being used, is commonly referred 
to as the \textit{Trusted Code Base (TCB)}. Multiple independent 
checkers can be employed to further increase confidence.

Relying on independent checkers as the TCB covers many of Lean's 
core applications, particularly those in which the proof itself 
is the main product:

\begin{itemize}
    \item Mathematical theorems.
    \item Verification of existing algorithms, protocols, or hardware.
    \item Verification of software written in other programming 
    languages and transformed into Lean statements using tools 
    such as \href{https://github.com/AeneasVerif/aeneas}{Aeneas}.
    The tool in question may introduce additional assumptions 
    concerning how the Lean statements relate to the original 
    software. These assumptions consequently become part of the TCB.
\end{itemize}

An extended use case of Lean, which also entails an extended TCB, 
arises from combining Lean's roles as a theorem prover and as a 
programming language. Examples include:

\begin{itemize}
    \item Executing verified Lean code as native binaries.
    \item Using the \texttt{native\_decide} tactic, which is tracked 
    as an additional axiom for transparency.
\end{itemize}

In these cases, the TCB is extended to include the Lean compiler 
(written in Lean), the Lean runtime (written in C++), and the 
code-generation backend (LLVM by default), as well as any other 
external code linked into the program.

Standard hardening techniques can be employed to increase 
confidence in these components, particularly when combining 
verified and unverified code. In this context, see also the 
\texttt{LEAN\_MI\_SECURE} CMake configuration option, which 
provides allocator hardening for Lean versions 4.31.0 and later.

\paragraph{How do I validate a Lean proof}
ref: \href{https://lean-lang.org/faq/#validation}{FAQ}.

The answer to that question depends on your threat model: Did you write the proof yourself, or someone you trust, or someone who is trying to trick you or AI-generated proofs or programs. Ref: \href{https://lean-lang.org/doc/reference/latest/find/?domain=Verso.Genre.Manual.section&name=validating-proofs}{FAQ}

Furthermore it is important to distinguish the question does the theorem have a valid proof from “what does the theorem statement mean. 

Let's look at our quicksort algorithm \ref{qsort_validation}. If one forgets to specify any sorting algorithm is a permutation of the input list, we get a proof for the isSorted criteris only, what is not sufficent for quicksort. Lean can not help here.  


